\chapter{Examples that distinguish tangents and supports}\label{sec:examples-rectifiable-tangents}
The tangent theorems describe almost-everywhere measure geometry, with precise normalization and hypotheses. The following examples explain why uniform jet conditions, local area bounds, positive weights and the mode of convergence matter. The framework is \cite{Sim18}, Chapter 2, Remark 1.7, pp.~47--48, and Chapter 3, pp.~85--92.

\begin{example}[Pointwise jet remainders do not suffice for a smooth extension]\label{ex:whitney-nonuniform-jet}
Let $C=\{0\}\cup\{1/k:k\geq1\}$, set $h(0)=0$, $h(1/k)=(-1)^k k^{-3/2}$ and $a=0$ on $C$. Both $h$ and $a$ are continuous. At every fixed $x\in C$ the first-order remainder tends to zero as $y\to x$ within $C$, but the jet has no $C^1(\mathbb R)$ extension. This is the example in \cite{Sim18}, Chapter 2, Remark 1.7(1), p.~47.
\end{example}
\begin{proof}
Every nonzero point of $C$ is isolated, so its pointwise condition is vacuous at sufficiently small distances. At zero the normalized remainder is $|h(1/k)|/(1/k)=k^{-1/2}\to0$. These values also prove continuity of $h$ at zero. However the absolute slope between adjacent points is
\[
 \frac{k^{-3/2}+(k+1)^{-3/2}}{1/k-1/(k+1)}
 =k(k+1)\bigl(k^{-3/2}+(k+1)^{-3/2}\bigr)\longrightarrow\infty.
\]
If a $C^1$ extension existed, the ordinary mean value theorem on each adjacent interval would produce a derivative of this absolute size at a point tending to zero. Continuity of the derivative near zero would make those values bounded, a contradiction. Equivalently the supremum in Definition \ref{def:whitney-first-jet} fails to tend to zero on the compact set $C$. The missing hypothesis is uniformity over nearby pairs, not continuity of the prescribed gradient alone.
\end{proof}

\begin{example}[A parabola converges weakly to its tangent but not in variation]\label{ex:parabola-weak-not-variation}
Let $M=\{(s,s^2):s\in\mathbb R\}$ and $\sigma_r=r^{-1}(\eta_{0,r})_\#(\mathcal H^1\llcorner M)$. Then $\sigma_r\rightharpoonup\mathcal H^1\llcorner P$ for $P=\{y=0\}$. Nevertheless, writing $|\alpha|$ for the total variation measure of a signed measure,
\[
 \left|\sigma_r-\mathcal H^1\llcorner P\right|(B_1(0))\longrightarrow4.
\]
Thus the tangent limit is not convergence in total variation, even for a smooth embedded curve.
\end{example}
\begin{proof}
The rescaled curve is $y=rs^2$, and the graph area formula gives $d\sigma_r=\sqrt{1+4r^2s^2}\,ds$ on it. Proposition \ref{prop:c1-graph-tangent-measure} gives weak convergence; its ball version gives $\sigma_r(B_1)\to2$. For every $r>0$, the parabola and $P$ meet only at zero, a point of zero length for both measures. Their restrictions to $B_1$ are mutually singular. Therefore the total variation of their difference on this ball is the sum of their masses, $\sigma_r(B_1)+\mathcal H^1(P\cap B_1)=\sigma_r(B_1)+2\to4$. This uses the full signed-measure variation convention, with no probability-theory factor of one half.
\end{proof}

\begin{example}[Weighted circles with infinite unweighted area near their limit]\label{ex:weighted-accumulating-circles}
In $\mathbb R^3$ let
\[
 C_j=\{(\cos t,\sin t,1/j):0\leq t<2\pi\},\qquad
 M=\bigcup_{j\geq1}C_j,\qquad\theta|_{C_j}=2^{-j}.
\]
The measure $\mu=\theta\mathcal H^1\llcorner M$ has total mass $2\pi$. The set $M$ is countably one-rectifiable but its unweighted length is infinite in every neighborhood of any point of the limiting circle $C_\infty=\{(\cos t,\sin t,0)\}$. Every point of $C_j$ has weighted tangent $2^{-j}\mathcal H^1$ on its tangent line. Every point of $C_\infty$ belongs to $\operatorname{supp}\mu$ but has one-dimensional density zero.
\end{example}
\begin{proof}
The circles are disjoint smooth curves, and summing $2\pi2^{-j}$ gives the total mass. At any point on a fixed $C_j$, its height is a positive distance from the heights of all other circles, so a sufficiently small neighborhood meets only this circle. Its weighted tangent follows from Proposition \ref{prop:c1-graph-tangent-measure} or the local smooth curve calculation.

By rotation take $q=(1,0,0)$ on the limiting circle. In each fixed ball about $q$, all sufficiently large $j$ contribute an arc of length bounded below by a positive constant depending only on that ball. There are infinitely many such circles, proving infinite unweighted length there. Every ball has positive weighted mass as well, so $q$ belongs to the support.

For $0<r<1/2$, only circles with $j>1/r$ can meet $B_r(q)$. The projection of a contributing arc onto its horizontal unit circle lies in the chordal ball of radius $r$ about $(1,0)$, whose arc length is $4\arcsin(r/2)\leq\pi r$. Thus
\[
 0\leq\frac{\mu(B_r(q))}{2r}
 \leq\frac\pi2\sum_{j>1/r}2^{-j}\longrightarrow0.
\]
The limiting circle is absent from $M$. This example illustrates both the usefulness of weighted tangents without locally finite unweighted area and the distinction between almost-everywhere density and density at every support point. No stationarity is asserted.
\end{proof}

\begin{example}[A finite rectifiable length measure can have a two-dimensional support]\label{ex:dense-circle-rectifiable-support}
There is a Borel countably one-rectifiable set $M\subset(0,1)^2$ with $0<\mathcal H^1(M)<\infty$ such that
\[
 \operatorname{supp}(\mathcal H^1\llcorner M)=[0,1]^2.
\]
Nevertheless $M$ has an approximate tangent line and unit one-dimensional density at $\mathcal H^1$-almost every point. Therefore an almost-everywhere tangent theorem is not a smoothness theorem for the topological support.
\end{example}
\begin{proof}
Enumerate the rational points $q_j$ of $(0,1)^2$ without repetitions. Choose $0<r_j<\min\{2^{-j},\tfrac12\operatorname{dist}(q_j,\partial(0,1)^2)\}$, and let $C_j$ be the circle of radius $r_j$ centered at $q_j$. Put $M=\bigcup_jC_j$. This is a countable union of compact smooth curves, contained in the open square, and $\mathcal H^1(M)\leq2\pi\sum_jr_j<\infty$. Its measure is positive because it contains a circle of positive length.

Every open ball contained in the square contains some whole circle: choose infinitely many rational centers in a smaller concentric ball, and then one with sufficiently large index that its radius bound $2^{-j}$ is smaller than the remaining distance to the ball's edge. Thus every such ball has positive restricted length. Every neighborhood of a boundary point of the closed square contains an interior ball of this type. Outside the closed square the measure vanishes on a neighborhood. These observations prove the exact support identity. Countable rectifiability and finite area permit Theorem \ref{thm:rectifiable-forward-tangents} and Corollary \ref{cor:rectifiable-local-unit-density} to be applied, proving the almost-everywhere conclusions despite the dense accumulation of other circles.
\end{proof}

\begin{example}[A zero weight can hide an unrectifiable component]\label{ex:zero-weight-hides-dust}
Let $P$ be the horizontal axis and let $K'=K+(0,2)$, where $K$ is the four-corner set of Definition \ref{def:four-corner-cantor}. Set $M=P\cup K'$, with weight $\theta=1$ on $P$ and $\theta=0$ on $K'$. Then $\mu=\theta\mathcal H^1\llcorner M$ is line measure on $P$, although $M$ is not countably one-rectifiable. If zero weights were allowed in the equality defining a weighted tangent, every line would satisfy that equality at every point of $K'$. Positivity in Definition \ref{def:weighted-approximate-tangent} is therefore necessary for both uniqueness and the characterization.
\end{example}
\begin{proof}
The weighted measure identity is immediate. The translated dust has finite positive length and is purely one-unrectifiable by Theorem \ref{thm:four-corner-pure-unrectifiability}; if $M$ were rectifiable, its measurable subset $K'$ would be rectifiable, contradicting that result. The dust is a positive distance from $P$. At a point of $K'$, a compactly supported scaled test therefore sees no weighted mass for all sufficiently small radii, so its limiting integral is zero. The putative right side $\theta(x)\int_L\phi$ is also zero for every line $L$. At a point of $P$ the usual line tangent holds. Thus simply removing the positive-weight hypothesis would make the purported tangent condition compatible with the unrectifiable dust. The dust result follows the example in \cite{Mor16}, pp.~36--37.
\end{proof}

\begin{example}[Changing one value of a weight changes a pointwise tangent formula]\label{ex:null-weight-value-tangent}
On an $m$-plane $P$, let $\theta=1$ except that $\theta(0)=7$, where $m\geq1$. The measure $\theta\mathcal H^m\llcorner P$ is still $\mathcal H^m\llcorner P$. Its normalized tangent at zero has coefficient one, not seven. The weighted tangent formula with the specified pointwise value $\theta(0)$ therefore fails there, while holding at every other point.
\end{example}
\begin{proof}
A singleton has zero $\mathcal H^m$ measure for $m\geq1$, so the measures are equal. Dilation centered at zero leaves normalized plane measure unchanged. For a nonnegative continuous compactly supported test with positive integral on $P$, multiplying that integral by seven gives a different value; no change of plane can repair the coefficient, as the unit-ball mass would still be seven times $\omega_m$. At a nonzero point the prescribed value is one and the same plane calculation applies. This is why theorems about a measurable weight and its density are almost-everywhere assertions and must not depend on arbitrarily reassigned exceptional values.
\end{proof}

\begin{example}[An artificial patch edge has half the plane density]\label{ex:halfplane-patch-edge}
Let $P$ be an $m$-plane, $m\geq1$, and let $e\in P$ be a unit vector. For $N=\{z\in P:z\cdot e\geq0\}$, every point of its edge $\{z\in P:z\cdot e=0\}$ has density $1/2$ with respect to $\mathcal H^m\llcorner N$, and its normalized tangent measure is a half-plane measure. It is not the full-plane tangent of Definition \ref{def:approximate-tangent}. The edge is $\mathcal H^m$-null.
\end{example}
\begin{proof}
Translation by an edge point and dilation preserve the half-plane and its normalized measure. Reflection in the edge exchanges the two halves of the plane ball and preserves Hausdorff measure, so its ball mass is exactly $\omega_mr^m/2$. The measure limit is therefore the same half-plane measure at every scale, whose ball mass excludes any unit full-plane limit. The edge is an $(m-1)$-plane and has zero $\mathcal H^m$ measure, directly from diameter covers on bounded pieces. This explains why the finite-area patch arguments use interior points of open base domains: an artificial truncation edge should not be treated as an interior smooth tangent point.
\end{proof}
